\documentclass[tikz,border=4pt]{standalone}
\usepackage{ctex}
\usepackage{amsmath}
\usetikzlibrary{calc, arrows.meta, decorations.pathreplacing}

\begin{document}
% part14/06 含绝对值数轴几何意义
% |x-1|+|x-3|=4 的临界点分段示意
\begin{tikzpicture}[scale=1.2, thick]
  % 数轴
  \draw[->] (-1.5,0) -- (6.5,0) node[right, font=\small] {};
  \foreach \x in {-1,0,1,2,3,4,5,6}{
    \draw (\x, 2pt) -- (\x, -2pt) node[below, font=\small] {$\x$};
  }

  % 临界点标记（1和3）
  \fill[red] (1,0) circle(2.5pt);
  \fill[red] (3,0) circle(2.5pt);
  \node[above, red, font=\small] at (1,0.08) {临界点 $1$};
  \node[above, red, font=\small] at (3,0.08) {临界点 $3$};

  % 三段分区标注
  \draw[cyan!60, thick, decorate, decoration={brace,amplitude=4pt}]
    (-1.3,0.25) -- (0.95,0.25) node[midway, above, cyan!60!black, font=\footnotesize, yshift=2pt] {情形1: $x<1$};
  \draw[orange!80, thick, decorate, decoration={brace,amplitude=4pt}]
    (1.05,0.25) -- (2.95,0.25) node[midway, above, orange!80!black, font=\footnotesize, yshift=2pt] {情形2: $1\leq x\leq3$};
  \draw[green!60!black, thick, decorate, decoration={brace,amplitude=4pt}]
    (3.05,0.25) -- (5.5,0.25) node[midway, above, green!60!black, font=\footnotesize, yshift=2pt] {情形3: $x>3$};

  % 解标注（x=0和x=4）
  \fill[blue] (0,0) circle(3pt);
  \node[below left, blue, font=\small] at (0,-0.15) {$x=0$ ✓};
  \fill[blue] (4,0) circle(3pt);
  \node[below right, blue, font=\small] at (4,-0.15) {$x=4$ ✓};

  % 距离标注
  \draw[<->, blue, thin] (0,-0.6) -- (4,-0.6) node[midway, below, font=\small, blue] {$|0-4|=4$};

  % 几何意义说明
  \node[font=\small] at (2.5,-1.2) {$|x-1|$ = $x$到$1$的距离，$|x-3|$ = $x$到$3$的距离};
  \node[font=\small] at (2.5,-1.6) {$|x-1|+|x-3|=4$ 的解：$x=0$ 或 $x=4$};

  % 中间段说明
  \node[orange!80!black, font=\footnotesize] at (2,-0.6) {$|x-1|+|x-3|=2$（恒为2$\neq 4$）};
\end{tikzpicture}
\end{document}
